\chapter{线性空间}
\section{线性空间}
\begin{definition}{线性空间}{linear space}
	数域$\FF$上的线性空间(linear space)~$V$是具有加法$+:V \times V \to V$和数乘$\cdot :\FF \times V \to V$运算且满足以下公理的集合。
	\begin{compactenum}
		\item 加法交换律\qqqquad\quad~$x + y = y + x;$
		\item 加法结合律\qqqquad\quad~$x + (y + z) = (x + y) + z;$
		\item 加法零元\qqqquad\qquad~$x + 0 = x;$
		\item 加法逆元\qqqquad\qquad~$x + (-x) = 0;$
		\item 数乘单位元\qqqquad\quad~$1x = x;$
		\item 数乘结合律\qqqquad\quad~$(c_1c_2)x = c_1(c_2x);$
		\item 数乘对向量的分配律\quad$c(x + y) = cx + cy;$
		\item 数乘对标量的分配律\quad$(c_1 + c_2)x = c_1x + c_2x.$
	\end{compactenum}
\end{definition}

\begin{example}
	{}{}
	$\FF^n$和$\FF^{m\times n}$都是线性空间。
\end{example}

\begin{definition}{子空间}{subspace}
	给定线性空间$V$的非空子集$V_\mathrm s \subset V$，若其对于加法和数乘封闭：$\forall v,w \in V_\mathrm s,\forall c \in \FF$
	\begin{equation}\label{eq:la-subspace-closure}
		v + w \in V_\mathrm s,\quad cv \in V_\mathrm s,
	\end{equation}
	则$V_\mathrm s$为$V$的子空间(subspace)。子空间中元素的线性组合都在同一个子空间，
\end{definition}

\begin{corollary}
	子空间必然包含零向量。因为若$v \in V_\mathrm s$，则$v + (-v) = 0 \in V_\mathrm s$。
\end{corollary}

\begin{remark}
	线性空间$V$的一个子集一般不是子空间，但我们可以通过其构造出子空间。
\end{remark}

\begin{definition}{线性扩张}{linear span}
	$S \subset V$的线性扩张(linear span)是$S$中向量的所有有限线性组合的集合，记作$\Span(S)$；约定$\Span(\varnothing) = \{0\}$。
\end{definition}

\begin{corollary}
	$V$的子集的线性扩展是$V$的子空间。
\end{corollary}

\section{线性独立、基和维度}
本章以下讨论有限维线性空间。
\begin{definition}{线性独立}{linear independent}
	$n$个向量$\{v_i\}$是线性独立的(linear independent)，当且仅当
	\begin{equation}\label{eq:la-linear-independence}
		\sum_{i=1}^nx_iv_i = 0,
	\end{equation}
	只在$x_i = 0$时成立，即只有零解。$n$个向量$\{v_i\}$不是线性独立，那么他们是线性相关的(linear correlate)。
	
	等价描述：集合中每一个向量都不能写成其它向量的线性组合。
\end{definition}
\begin{remark}
	向量是否线性独立同数域的选择有关。
\end{remark}
\begin{definition}{线性空间的基}{base}
	线性空间$V$的基(base)是一组线性无关的向量$\{v_i\}$，并且他们张成整个线性空间$V$。
\end{definition}
\begin{example}
	{}{}
	$\{e_1,\ldots,e_n\}$构成$\RR^n$的一组基。
\end{example}
\begin{definition}{线性空间的维度}{dimension}
	线性空间的维度(dimension)是一组基中向量的个数，记作$\dim(V)$。
\end{definition}
\begin{theorem}{维度的确定性}{certainty of dimension}
	线性空间的维度和基的选取无关。
\end{theorem}
\begin{proof}
	若线性空间$V$存在两组基$\{v_1,\ldots,v_m\},\{w_1,\ldots,w_n\}$元素个数不等，不妨设$m > n$。
	因为$\{w_i\}$是基，$\{v_i\}$可以被表示为其线性组合
	\[
		v_i = \sum_{j=1}^nw_ja_{ji},\quad\forall i.
	\]
	考虑线性组合
	\[
		\sum_{i=1}^mx_iv_i = \sum_{i=1}^m\sum_{j=1}^nx_iw_ja_{ji} = 0.
	\]
	因为$\{w_i\}$线性无关，故
	\[
		\sum_{i=1}^ma_{ji}x_i = 0,\quad\forall j,
	\]
	但是其未知数的个数$m >$方程的个数$n$，系数矩阵$A$一定有自由列，所以$x$有非零解。这与$\{v_i\}$线性无关矛盾！故$m = n.$
\end{proof}

\begin{theorem}{}{}
	设$V_1,V_2$是同一有限维线性空间的子空间。若$\forall v \in V_1$可以写成$V_2$中向量的线性组合，则$V_1 \subseteq V_2$，从而$\dim(V_1) \leqslant \dim(V_2)$。
\end{theorem}

\begin{proof}
	这个定理是trivial的，证明留给读者。
\end{proof}

\begin{definition}{变换矩阵}{transfomation matrix}
	不同基之间的变换相应的矩阵称为变换矩阵(transfomation matrix)。
\end{definition}

\section{矩阵\textit{A}的四个子空间}

给定矩阵$A \in \FF^{m\times n}$，可由其得到四个子空间：列空间、行空间、零空间和左零空间。
\begin{definition}{列空间和行空间}{column space}
	矩阵$A$的列空间(column space)是$A$的所有列的线性组合的集合，记作$\colsp(A)$：
	\begin{equation}
		\colsp(A) \equiv \set{Ax}{x \in \FF^n}.
	\end{equation}
	\tcblower
	矩阵$A$的行空间(row space)是$A$的所有行的线性组合的集合，%由于转置并不影响性质，行空间
	可用$\colsp(A\transp)$表示。
\end{definition}

\begin{corollary}
	$\colsp(A)$是$\FF^m$的子空间，$\colsp(A\transp)$是$\FF^n$的子空间。
\end{corollary}

\begin{remark}
	线性方程组$Ax = b$有解$\iff b \in \colsp(A)$。
\end{remark}

\begin{definition}{零空间和左零空间}{null space}
	矩阵$A$的零空间(null space)是所有满足$Ax = 0$的$x$的集合，记作$\nullsp(A)$：
	\begin{equation}
		\nullsp(A) \equiv \set x{Ax = 0}.
	\end{equation}
	\tcblower
	矩阵$A$的左零空间(left null space)是所有满足$x\transp A = 0$的$x$的集合，可用$\nullsp(A\transp)$表示。
\end{definition}

\begin{corollary}
	$\nullsp(A)$是$\FF^n$的子空间，$\nullsp(A\transp)$是$\FF^m$的子空间。
\end{corollary}

\begin{example}{子空间的基}{}
	$\nullsp(A)$的基：显然$\nullsp(A) = \nullsp(\rref(A))$，而$\rref(A)$的每个自由列各给出一个$\nullsp(\rref(A))$的基向量。若第$j$列为自由列，则其对应的基$x$的分量为：
	\begin{equation}\label{eq:la-null-basis-components}
		x_i = \begin{cases}
			\delta_{ij},&\text{第$i$列为自由列}\\
			-\rref(A)_{kj},&\text{第$i$列为主列，主元位于第$k$行}
		\end{cases}
	\end{equation}
	%$\dim(\rref(A))=\rref(A)$中自由列的数量；
	\tcblower
	$\colsp(A)$的基：先找出$\rref(A)$的主元列编号，再取原矩阵$A$中相应编号的列。行变换一般改变列空间，不能直接取$\rref(A)$的列作为$\colsp(A)$的基。
\end{example}

\begin{example}
	{用矩阵配平化学方程式}{}
	考虑下列反应式的配平问题，这里只讨论元素守恒：
	\[
		\mathrm{P_4+CuSO_4+H_2O\to Cu_3P+H_3PO_4+H_2SO_4}.
	\]
	令反应物的带符号系数$\nu$为正、生成物的为负。由反应前后$\mathrm{H,O,P,S,Cu}$元素守恒可得：
	\[
		\begin{bmatrix}
			&&2&&3&2\\
			&4&1&&4&4\\
			4&&&1&1\\
			&1&&&&1\\
			&1&&3
		\end{bmatrix}
		\begin{bmatrix}
			\nu_\mathrm{P_4}\\\nu_\mathrm{CuSO_4}\\\nu_\mathrm{H_2O}\\\nu_\mathrm{Cu_3P}\\\nu_\mathrm{H_3PO_4}\\\nu_\mathrm{H_2SO_4}
		\end{bmatrix}
		=
		\begin{bmatrix}
			0\\0\\0\\0\\0
		\end{bmatrix}
		\implies
		\begin{bmatrix}
			\nu_\mathrm{P_4}\\\nu_\mathrm{CuSO_4}\\\nu_\mathrm{H_2O}\\\nu_\mathrm{Cu_3P}\\\nu_\mathrm{H_3PO_4}\\\nu_\mathrm{H_2SO_4}
		\end{bmatrix}
		= k\begin{bmatrix}
			11\\60\\96\\-20\\-24\\-60
		\end{bmatrix}.
	\]
	这相当于求系数矩阵的零空间的基，得到配平的化学方程式为：
	\[
		\mathrm{11\,P_4+60\,CuSO_4+96\,H_2O=20\,Cu_3P+24\,H_3PO_4+60\,H_2SO_4}.
	\]
	\tcblower
	但是这种方法也不是万能的，比如高锰酸钾与双氧水的反应：
	\[
		\mathrm{MnO_4^-+H_2O_2+H^+\to Mn^{2+}+H_2O+O_2\uparrow}.
	\]
	由反应前后$\mathrm{H,O,Mn}$元素守恒及电荷守恒可得：
	\[
		\begin{bmatrix}
			&2&1&&2\\
			4&2&&&1&2\\
			1&&&1\\
			-1&&1&2
		\end{bmatrix}
		\begin{bmatrix}
			\nu_\mathrm{MnO_4^-}\\\nu_\mathrm{H_2O_2}\\\nu_\mathrm{H^+}\\\nu_\mathrm{Mn^{2+}}\\\nu_\mathrm{H_2O}\\\nu_\mathrm{O_2}
		\end{bmatrix} = \begin{bmatrix}
			0\\0\\0\\0
		\end{bmatrix}
		\implies
		\begin{bmatrix}
			\nu_\mathrm{MnO_4^-}\\\nu_\mathrm{H_2O_2}\\\nu_\mathrm{H^+}\\\nu_\mathrm{Mn^{2+}}\\\nu_\mathrm{H_2O}\\\nu_\mathrm{O_2}
		\end{bmatrix} = k_1\begin{bmatrix}
			2\\5\\6\\-2\\-8\\-5
		\end{bmatrix} + k_2\begin{bmatrix}
			0\\2\\0\\0\\-2\\-1
		\end{bmatrix}.
	\]
	零空间为二维，单靠这些守恒方程不能唯一确定反应的系数比，这是因为该方程可以线性叠加双氧水分解反应($\mathrm{2\,H_2O_2=2\,H_2O+O_2\uparrow}$)。
	若进一步指定氧原子的来源，假定$\mathrm{H_2O_2}$中的氧全部进入$\mathrm{O_2}$，而$\mathrm{MnO_4^-}$中的氧全部进入水，可用同位素标签表示为：
	\[
		\mathrm{Mn{}^{16}O_4^-+H_2{}^{18}O_2+H^+\to Mn^{2+}+H_2{}^{16}O+{}^{18}O_2\uparrow}.
	\]	
	便可由$\mathrm{{}^{16}O,{}^{18}O}$分别守恒确定$\nu_\mathrm{MnO_4^-}:\nu_\mathrm{H_2O_2} = 2:5$。
	\[
		\mathrm{2\,MnO_4^-+5\,H_2O_2+6\,H^+=2\,Mn^{2+}+8\,H_2O+5\,O_2\uparrow}.
	\]
\end{example}


\section{矩阵的秩、线性代数基本定理}

\begin{definition}{矩阵的秩}{rank}
	矩阵$A$的秩(rank)定义为列空间的维数：%\footnote{后面很快会证明行空间列空间维数相等。}：
	\begin{equation}\label{eq:la-matrix-rank}
		\rank(A) := \dim(\colsp(A)).
	\end{equation}
	若矩阵的秩等于行数，则称其行满秩(full row rank)；
	若秩等于列数，则称其列满秩(full column rank)。
\end{definition}
\begin{theorem}
	{线性方程组$Ax = b$的完整解}{}
	若线性方程组$Ax = b$有解，其通解可写为一个特解$x_{\mathrm p}$与齐次方程解$x_{\mathrm n}$的和：
	\begin{equation}\label{eq:la-inhomogeneous-solutions}
		x = x_{\mathrm p} + x_{\mathrm n},
	\end{equation}
	%其中$x_{\mathrm n}\in\nullsp(A)$，
	特解满足$Ax_{\mathrm p} = b$，其各个分量可以由约化的增广矩阵$[\rref(A)\enspace b']$得到：
	\begin{equation}\label{eq:la-particular-solution-components}
		(x_{\mathrm p})_i = \begin{cases}
			0,&\text{第$i$列为自由列}\\
			b'_j,&\text{第$i$列为主列，其主元1在第$j$行}
		\end{cases}
	\end{equation}
	而齐次解满足$Ax_{\mathrm n} = 0$，可以写成零空间的基的线性组合。
\end{theorem}
\begin{corollary}
	~
	\begin{itemize}
		\item 若$A$列满秩，则$\rref(A)$没有自由列，$\nullsp(A) = \{0\}$。当$b \in \colsp(A)$时有唯一解，否则无解；
		\item 若$A$行满秩，则$\rref(A)$没有零行，$\colsp(A) = \FF^m$，$\forall b$都有解，有唯一解或无穷多解。
	\end{itemize}
\end{corollary}
\paragraph{四个子空间的维度}已经知道，初等行变换就是用初等矩阵$E$左乘$A$，相应的，列变换就是$E$右乘$A$。
\begin{theorem}{初等变换和子空间}{}
	\begin{compactitem}
		\item $\nullsp(A) = \nullsp(EA)$，因为
		\[
			Ax = 0 \iff EAx = 0.
		\]
		\item $\dim(\colsp(A)) = \dim(\colsp(EA))$，因为
		\[
			\{v_i\}\text{~是~}\colsp(A)\text{~一组基} \iff \{Ev_i\}\text{~是~}\colsp(EA)\text{~一组基}.
		\]
		\item $\colsp(A) = \colsp(AE)$，因为
		\[
			AE\text{~的每一列} \in \colsp(A)\text{~且~}A = (AE)E\inv\text{~的每一列} \in \colsp(AE)
		\]
		\item $\dim(\nullsp(A)) = \dim(\nullsp(AE))$，因为可由$Ax = 0 \iff AE(E\inv x) = 0$推出
		\[
			\{v_i\}\text{~是~}\nullsp(A)\text{~一组基} \iff \{E\inv v_i\}\text{~是~}\nullsp(AE)\text{~一组基}
		\]
	\end{compactitem}
\end{theorem}
\begin{corollary}
	矩阵$A$在初等变换下，$\dim(\colsp(A))$和$\dim(\nullsp(A))$均不变，而行变换下$\nullsp(A)$不变，列变换下$\colsp(A)$不变。
	因此，可以将$A$先由行变换为$\rref(A)$，再列变换为
	\[
		\tilde I = \begin{bmatrix}
			I_r&0\\0&0
		\end{bmatrix} \in \FF^{m\times n}.
	\]
	显然，$\tilde I$的行秩$\dim(\colsp(\tilde I\transp)) =$列秩$\dim(\colsp(\tilde I)) = r$，且
	\begin{align*}
		\dim(\colsp(\tilde I\transp)) + \dim(\nullsp(\tilde I)) &= n;\\
		\dim(\colsp(\tilde I)) + \dim(\nullsp(\tilde I\transp)) &= m.
	\end{align*}
\end{corollary}
% 以上的这些量在初等变化下都不变，故
\begin{theorem}{线性代数基本定理·一}{fundamental theorems of linear algebra I}
	\begin{compactenum}
		\item 行秩=列秩：$\rank(A) = \dim(\colsp(A)) = \dim(\colsp(A\transp));$
		\item $\dim(\colsp(A\transp)) + \dim(\nullsp(A)) = n;$
		\item $\dim(\colsp(A)) + \dim(\nullsp(A\transp)) = m.$
	\end{compactenum}
\end{theorem}
\begin{corollary}
	方阵$A$可逆$\iff A$满秩。
\end{corollary}
\begin{notecaptionblock}
    \tikzfigure{four-subspaces}
    \captionof{figure}{$A$的四个子空间}
\end{notecaptionblock}

\subsectionstar{矩阵的秩的不等式}

% 参考：\url{https://zhuanlan.zhihu.com/p/55206421}

\begin{theorem}{}{}
	对于同型矩阵$A,B$有
	\begin{equation}
		\rank(A + B) \leq \rank(A) + \rank(B).
	\end{equation}
\end{theorem}
\begin{proof}
	设$\colsp(A),\colsp(B)$的一组基分别为$\{a_1,\ldots,a_{\rank(A)}\}$和$\{b_1,\ldots,b_{\rank(B)}\}$，则$A + B$中的列必然可以由$\{a_1,\ldots,a_{\rank(A)},b_1,\ldots,b_{\rank(B)}\}$的线性组合表示出，故不等式成立。
\end{proof}
\begin{theorem}{Sylvester不等式}{Sylvester inequality}
	对$A \in \FF^{m\times n},B \in \FF^{n\times p}$有
	\begin{equation}
		\label{eqn:Sylvester inequality}
		\rank(AB) \geq \rank(A) + \rank(B) - n.
	\end{equation}
\end{theorem}
\begin{proof}
	注意到行/列等价关系：
	\[
		\begin{bmatrix}
			I_n\\ &AB
		\end{bmatrix} \sim \begin{bmatrix}
			I_n\\A&AB
		\end{bmatrix} \sim \begin{bmatrix}
			I_n&-B\\A
		\end{bmatrix}.
	\]
	两边取秩即得
	\[
		\rank(AB) + n = \rank\biggp{\begin{bmatrix}
			I_n&-B\\A
		\end{bmatrix}} \geq \rank(A) + \rank(B).
		\qedhere
	\]
\end{proof}

\begin{theorem}{Frobenius不等式}{Frobenius inequality}
	对于$A \in \FF^{m\times n},B \in \FF^{n\times p},C \in \FF^{p\times q}$，有
	\begin{equation}
		\label{eqn:Frobenius inequality}
		\rank(ABC) \geq \rank(AB) + \rank(BC) - \rank(B).
	\end{equation}
\end{theorem}
\begin{proof}
	注意到行/列等价关系：
	\[
		\begin{bmatrix}
			ABC\\ &B
		\end{bmatrix} \sim \begin{bmatrix}
			ABC&AB\\ &B
		\end{bmatrix} \sim \begin{bmatrix}
			&AB\\-BC&B
		\end{bmatrix},
	\]
	两边取秩即得
	\[
		\rank(ABC) + \rank(B) = \rank\biggp{\begin{bmatrix}
			&AB\\-BC&B
		\end{bmatrix}} \geq \rank(AB) + \rank(BC).
		\qedhere
	\]
\end{proof}

\begin{corollary}
	对Frobenius不等式\eqref{eqn:Frobenius inequality}取$B = I_n$即可得到Sylvester不等式\eqref{eqn:Sylvester inequality}。
\end{corollary}

